\chapter{4.10 关于开展供应链创新与应用试点的通知(商务部等)} %2018

\begin{center}
商务部市场体系建设司 \\
商建函[2018]142号
\end{center}

\noindent 各省、自治区、直辖市、计划单列市及新疆生产建设兵团商务、工业和信息化、环境保护、农业、质量技术监督（市场监管）部门，中国人民银行各分行、营业管理部、各省会（首府）城市中心支行、各副省级城市中心支行，各银监局，中国物流与采购联合会各分支机构：

根据《国务院办公厅关于积极推进供应链创新与应用的指导意见》（国办发〔2017〕84号）要求，商务部、工业和信息化部、生态环境部、农业农村部、人民银行、国家市场监督管理总局、中国银行保险监督管理委员会和中国物流与采购联合会决定开展供应链创新与应用试点。现将有关事项通知如下：

\section{总体要求}

  （一）指导思想。

全面贯彻党的十九大精神，以习近平新时代中国特色社会主义思想为指导，落实国务院关于推进供应链创新与应用的决策部署，以供给侧结构性改革为主线，完善产业供应链体系，高效整合各类资源和要素，提高企业、产业和区域间的协同发展能力，适应引领消费升级，激发实体经济活力，在现代供应链领域培育新增长点、形成新动能，助力建设现代化经济体系，推动经济高质量发展。

  （二）总体思路。

试点包括城市试点和企业试点，试点实施期为2年。试点城市的主要任务是出台支持供应链创新发展的政策措施，优化公共服务，营造良好环境，推动完善产业供应链体系，并探索跨部门、跨区域的供应链治理新模式。试点企业的主要任务是应用现代信息技术，创新供应链技术和模式，构建和优化产业协同平台，提升产业集成和协同水平，带动上下游企业形成完整高效、节能环保的产业供应链，推动企业降本增效、绿色发展和产业转型升级。

通过城市试点和企业试点，在若干关系国计民生、消费升级和战略新兴的重点产业，推动形成创新引领、协同发展、产融结合、供需匹配、优质高效、绿色低碳、全球布局的产业供应链体系，促进发展实体经济，助力供给侧结构性改革，筑牢现代化经济体系的坚实基础。

{ \bf （三）试点目标。 }

通过试点，打造“五个一批”，即创新一批适合我国国情的供应链技术和模式，构建一批整合能力强、协同效率高的供应链平台，培育一批行业带动能力强的供应链领先企业，形成一批供应链体系完整、国际竞争力强的产业集群，总结一批可复制推广的供应链创新发展和政府治理实践经验。

通过试点，现代供应链成为培育新增长点、形成新动能的重要领域，成为供给侧结构性改革的重要抓手，成为“一带一路”建设和形成全面开放新格局的重要载体。

\section{试点城市重点任务}

{ \bf （一）推动完善重点产业供应链体系。 }

一是建立健全农业供应链。结合本地特色农业，优先选择粮食、果蔬、茶叶、药材、乳制品、蛋品、肉品、水产品、酒等重要产品，立足区域特色优势，充分发挥农业产业化龙头企业示范引领作用，推动供应链资源集聚和共享，打造联结农户、新型农业经营主体、农产品加工流通企业和最终消费者的紧密型农产品供应链，构建完善全产业链各环节相互衔接配套的绿色可追溯农业供应链体系。

二是积极发展工业供应链。结合本地主导产业，优先选择钢铁、煤炭、水泥、玻璃等相关产业，推动企业打造供需对接、资源整合的供应链协同平台，提高产业协同效率，推动降成本、去库存和去产能，助力供给侧结构性改革。

在与消费升级密切相关的产业中，优先选择家电、汽车、电子、纺织等，推动企业构建对接个性化需求和柔性化生产的智能制造供应链协同平台，提高产品和服务质量，满足人民日益增长的美好生活需要。

针对必须抢占制高点的战略新兴产业，充分调动各方资源，打造合作紧密、分工明确、集成联动的政产学研一体化的供应链创新网络，推进大型飞机、机器人、发动机、集成电路等关键技术攻关和产业发展。

三是创新发展流通供应链。推动企业与供应商、生产商实现系统对接，构建流通与生产深度融合的供应链协同平台，实现供应链需求、库存和物流实时共享可视。

推动企业建设运营规范的商品现货交易平台，提供供应链增值服务，提高资源配置效率。促进传统实体商品交易市场转型升级，打造线上线下融合的供应链交易平台，促进市场与产业融合发展。

鼓励传统流通企业向供应链服务企业转型，建设供应链综合服务平台，提供研发、设计、采购、生产、物流和分销等一体化供应链服务，提高流通效率，降低流通成本。

推进城市居民生活供应链体系建设，发展集信息推送、消费互动、物流配送等功能为一体的社区商业，满足社区居民升级消费需求，提高居民生活智能化和便利化水平。

{ \bf （二）规范发展供应链金融服务实体经济。 }

推动供应链核心企业与商业银行、相关企业等开展合作，创新供应链金融服务模式，发挥上海票据交易所、中征应收账款融资服务平台和动产融资统一登记公示系统等金融基础设施作用，在有效防范风险的基础上，积极稳妥开展供应链金融业务，为资金进入实体经济提供安全通道，为符合条件的中小微企业提供成本相对较低、高效快捷的金融服务。

推动政府、银行与核心企业加强系统互联互通和数据共享，加强供应链金融监管，打击融资性贸易、恶意重复抵质押、恶意转让质物等违法行为，建立失信企业惩戒机制，推动供应链金融市场规范运行，确保资金流向实体经济。

{ \bf （三）融入全球供应链打造“走出去”战略升级版。 }

推动本地优势产业对接并融入全球供应链体系，开展更大范围、更高水平、更深层次的国际合作，向全球价值链中高端跃升，打造更具全球竞争力的产业集群。

支持和鼓励企业积极开展对外贸易与投资合作，加强与“一带一路”沿线国家的互联互通，设立境外研发中心、分销服务网络、物流配送中心、海外仓等，提高全球范围内供应链协同和配置资源的能力，促进重要资源能源、重要农产品、关键零部件来源的多元化和目标市场的多样化。

{ \bf （四）发展全过程全环节的绿色供应链体系。 }

推动深化政府绿色采购，行政机关和使用财政资金的其他组织应当优先采购和使用节能、节水、节材等环保产品、设备和设施，并建立相应的考核体系。研究制定重点产业企业绿色供应链构建指南，建立健全环保信用评价、信息强制性披露等制度，依法依规公开供应链全环节的环境违法信息。

支持环境保护技术装备、资源综合利用和环境服务等环境保护产业的发展。加大对绿色产品、绿色包装的宣传力度，鼓励开展“快递业+回收业”定向合作，引导崇尚自然、追求健康的消费理念，培育绿色消费市场。

{ \bf （五）构建优质高效的供应链质量促进体系。 }

加强供应链质量标准体系建设，推广《服务质量信息公开规范》和《服务质量评价工作通用指南》，建立供应链服务质量信息清单制度。加强全链条质量监管，开发适应供应链管理需求的质量管理工具，引入第三方质量治理机制，探索建立供应链服务质量监测体系并实施有针对性的质量改进。

引导企业树立质量第一的意识，提高服务质量，创新服务模式，优化服务流程，为客户提供安全、诚信、优质、高效的服务。鼓励企业加强供应链品牌建设，创建一批高价值供应链品牌。

{ \bf （六）探索供应链政府公共服务和治理新模式。 }

以完善政策、优化服务、加强监管为重点，改革相关体制机制，出台相关支持政策措施，宣传推广供应链思维、理念和技术，营造供应链创新与应用的良好环境。积极打造供应链公共服务平台，研究设立供应链创新产业投资基金，建设供应链科技创新中心，支持供应链前沿技术、基础软件、先进模式等的研究与推广。加强供应链标准体系、信用体系和人才体系等支撑建设。

\section{试点企业重点任务}

试点企业应围绕试点目标，发挥龙头带动作用，加强与供应链上下游企业的协同和整合，着力完善产业供应链体系，促进产业降本增效、节能环保、绿色发展和创新转型。

{ \bf （一）提高供应链管理和协同水平。 }

普及供应链思维，完善供应链管理制度，加强企业信息化升级，加强标准化建设，培养供应链专业人才，提高与上下游企业协同能力，形成分工协作的网络体系，积极“走出去”开展对外贸易投资合作，构建全球供应链，提升全球资源配置效率。

{ \bf （二）加强供应链技术和模式创新。 }

积极与高校、研究机构等开展合作，建设供应链研究中心或实验室，开展供应链技术创新和软硬件研发，推广应用供应链新技术、新模式，促进整个产业供应链数字化、智能化和国际化。

{ \bf （三）建设和完善各类供应链平台。 }

以平台为重要载体完善供应链体系，加强与上下游企业实现系统和数据对接，充分发挥供应链平台的资源集聚、供需对接和信息服务等功能，构建跨界融合的产业供应链生态。

{ \bf （四）规范开展供应链金融业务。 }

有条件的企业可加强与商业银行、平台企业等合作，创新供应链金融业务模式，优化供应链资金流，积极稳妥、依法依规开展供应链金融业务。

{ \bf （五）积极倡导供应链全程绿色化。 }

以全过程、全链条、全环节的绿色发展为导向，优先采购和使用节能、节水、节材等环保产品、设备和设施，促进形成科技含量高、资源消耗低、环境污染少的产业供应链。

\section{组织实施程序}

 （一）积极组织申报。

省级商务主管部门会同工业和信息化、环境保护、农业主管部门、人民银行、质量技术监督（市场监管）部门和原银监会派出机构，建立试点工作协调机制，共同组织本地各城市和企业申报，并将《供应链创新与应用城市试点方案》（附件1）、《供应链创新与应用试点企业申报表》（附件2）以及《申报供应链创新与应用试点单位汇总表》（附件3）（含电子版）一式八份报商务部（市场建设司）。申报截止日期为2018年 5月31日。

申报城市应拥有较好的产业基础，重点产业集群在全国具有较强的影响和带动能力，拥有比较完善的产业配套体系；具有较好的供应链发展软硬基础设施，现代信息技术应用规模较大。自由贸易试验区可独立申报。

申报企业应具有独立法人资格，具有较高的供应链管理能力，较完善的供应链管理制度，较强的供应链人才力量，建有业内较大影响的供应链平台，对行业发展具有重要影响和示范带动作用。中央企业可由国务院相关部门推荐申报，也可直接向商务部申报。

 （二）编制试点方案。

申报试点的城市应编制《供应链创新与应用城市试点方案》。实施方案要依托本市优势和特色产业，聚焦一项或多项试点任务，科学谋划供应链推动产业发展的总体思路、目标任务、试点内容和保障措施。

申报试点的企业应按要求填写《供应链创新与应用试点企业申报表》。申报企业要围绕试点目标和任务，提出加强供应链协同和整合、完善产业供应链体系的目标和举措。

{ \bf （三）确定试点城市和试点企业。 }

商务部会同工业和信息化部、生态环境部、农业农村部、人民银行、国家市场监督管理总局、中国银行保险监督管理委员会和中国物流与采购联合会共同组织专家对申报材料进行评审，通过竞争性择优确定试点城市和试点企业。

{ \bf （四）试点实施并及时报送进展情况。 }

各试点城市和试点企业应每季度上报试点情况。试点取得的重大进展，或遇到的重大问题和困难，应及时报告。省级商务主管部门将有关情况统一报商务部。试点企业是中央企业的，由中央企业直接上报商务部。

{ \bf （五）绩效评估和经验总结。 }

商务部将会同有关部门对试点进行中期和终期评估，对于中期评估表现优秀的城市和企业，研究给予相关激励政策；根据终期评估结果，对试点效果显著、绩效评估为优的城市授予“全国供应链创新与应用示范城市”称号，对行业引领能力强、绩效评估为优的企业授予“全国供应链创新与应用示范企业”称号。同时总结可复制推广的实践经验。

\section{工作要求}

{ \bf （一）加强组织领导和工作协调。 }

要加强试点的组织领导，建立相关部门参与的工作协调机制，制定试点方案，明确责任分工，加强统筹协调，确保工作顺利推进。各级商务、工业和信息化、环境保护、农业、人民银行、质量技术监督（市场监管）等部门和原银监会派出机构要密切配合，共同做好试点的组织实施、监督、评估和总结推广等工作，及时帮助协调解决试点的困难和问题。

 （二）加强业务指导和政策支持。

鼓励试点城市整合国家、省级预算内投资等各项资金，引导社会资本设立供应链创新产业投资基金，支持供应链创新和应用项目。试点城市有关部门在安排相关投资时，对于符合条件的供应链创新发展项目予以倾斜。

各级商务、工业和信息化、农业、环境保护和质量技术监督（市场监管）部门要加强对流通、工业、农业供应链以及供应链全程绿色化和供应链质量标准等方面的业务指导，研究出台相关财政、税收、金融、土地等方面的政策措施，加大对试点的支持力度。

各级人民银行和原银监会派出机构要加强对供应链金融发展的指导和监督，研究相关支持政策，对内外资企业和机构一视同仁，营造公平竞争的市场环境。优先支持符合条件的试点企业发行公司信用类债券，提高事中事后风险管理水平，推动供应链金融健康稳定发展。

各级商务等行业主管部门要指导和支持相关行业组织探索建立供应链绩效指数评价体系，为各地建立供应链公共服务平台提供技术支持，为试点城市和试点企业进行专业指导和业务培训等。

{ \bf （三）做好宣传和总结推广工作。 }

各级商务、工业和信息化、环境保护、农业、人民银行、质量技术监督（市场监管）部门和原银监会派出机构要及时梳理总结试点中出现的典型案例，在各类媒体进行宣传报道。试点中形成的先进模式和经验要及时推广，扩大试点效果。

\section{联系方式}
{ \bf （略） }
